\documentclass[12pt, a4paper]{article}
% 引入中文支持包
\usepackage{ctex} 

\usepackage[utf8]{inputenc}
\usepackage[T1]{fontenc}
\usepackage{amsmath, amssymb, amsthm}
\usepackage{geometry}
\geometry{left=2.5cm, right=2.5cm, top=2.5cm, bottom=2.5cm}

% 定义命令以便简化书写
\newcommand{\Rn}{\mathbb{R}^{n \times n}}
\newcommand{\R}{\mathbb{R}}

\begin{document}
	
	\section*{作业手稿整理}
	
	\begin{enumerate}
		\item \textbf{题目：} 对 $P \in \Rn$，若它满足 $P^2 = P$，称其为投影矩阵。证明：
		\begin{enumerate}
			\item[(1)] $I-P$ 为投影矩阵；
			\item[(2)] $R(I-P) = N(P)$，且 $N(I-P) = R(P)$；
			\item[(3)] $R(P) \cap N(P) = \{0\}$。
		\end{enumerate}
		
		\textbf{证明：}
		
		\begin{enumerate}
			\item[(1)] 
			\begin{align*}
				(I-P)^2 &= (I-P)(I-P) \\
				&= I^2 - IP - PI + P^2 \\
				&= I - 2P + P \quad (\text{因 } P^2=P) \\
				&= I - P
			\end{align*}
			故 $I-P$ 满足幂等性，为投影矩阵。
			
			\item[(2)] \textbf{证明 $R(I-P) = N(P)$：}
			\begin{itemize}
				\item 设 $y \in R(I-P)$，则 $\exists x \in \R^n$，使得 $y = (I-P)x$。
				\[
				Py = P(I-P)x = (P - P^2)x = (P-P)x = 0
				\]
				故 $y \in N(P)$，即 $R(I-P) \subseteq N(P)$。
				
				\item 设 $x \in N(P)$，即 $Px = 0$。
				\[
				x = Ix - Px + Px = (I-P)x + Px
				\]
				由于 $Px=0$，故 $x = (I-P)x$。这意味着 $x$ 可以表示为 $(I-P)$ 的像，即 $x \in R(I-P)$。
				故 $N(P) \subseteq R(I-P)$。
			\end{itemize}
			综上，$R(I-P) = N(P)$。
			
			\textbf{证明 $N(I-P) = R(P)$：}
			\begin{itemize}
				\item 设 $y \in R(P)$，则 $\exists x \in \R^n$，使得 $Px = y$。
				\[
				(I-P)y = (I-P)Px = (P - P^2)x = 0
				\]
				故 $y \in N(I-P)$，即 $R(P) \subseteq N(I-P)$。
				
				\item 设 $x \in N(I-P)$，则 $(I-P)x = 0 \Rightarrow x - Px = 0 \Rightarrow x = Px$。
				这表明 $x$ 在 $P$ 的像空间中（$x = P(x)$），即 $x \in R(P)$。
				故 $N(I-P) \subseteq R(P)$。
			\end{itemize}
			综上，$N(I-P) = R(P)$。
			
			\item[(3)] \textbf{证明 $R(P) \cap N(P) = \{0\}$：}
			设 $x \in R(P) \cap N(P)$。
			\begin{itemize}
				\item 由 $x \in N(P)$ 知，$Px = 0$。
				\item 由 $x \in R(P)$ 知，$\exists y \in \R^n$，使得 $x = Py$。
			\end{itemize}
			则：
			\[
			Px = P(Py) = P^2y = Py = x
			\]
			结合 $Px=0$，可得 $x = 0$。
			故 $R(P) \cap N(P) = \{0\}$。
		\end{enumerate}
		
		\item \textbf{题目：} 计算最小二乘问题 $\min \|Ax - b\|_2^2$ 的解。
		\[
		A = \begin{pmatrix} -1 & -1 & 1 \\ 1 & 3 & 3 \\ -1 & -1 & 5 \\ 1 & 3 & 7 \end{pmatrix}, \quad 
		b = \begin{pmatrix} 1 \\ 1 \\ 1 \\ 1 \end{pmatrix}
		\]
		
		\textbf{解：} 
		易见 $R(A)=3$（列满秩），最小二乘解满足法方程 $A^T A x = A^T b$。
		
		计算 $A^T A$：
		\[
		A^T A = \begin{pmatrix} -1 & 1 & -1 & 1 \\ -1 & 3 & -1 & 3 \\ 1 & 3 & 5 & 7 \end{pmatrix}
		\begin{pmatrix} -1 & -1 & 1 \\ 1 & 3 & 3 \\ -1 & -1 & 5 \\ 1 & 3 & 7 \end{pmatrix}
		= \begin{pmatrix} 4 & 8 & 4 \\ 8 & 20 & 24 \\ 4 & 24 & 84 \end{pmatrix}
		\]
		
		计算 $A^T b$：
		\[
		A^T b = \begin{pmatrix} -1 & 1 & -1 & 1 \\ -1 & 3 & -1 & 3 \\ 1 & 3 & 5 & 7 \end{pmatrix}
		\begin{pmatrix} 1 \\ 1 \\ 1 \\ 1 \end{pmatrix}
		= \begin{pmatrix} 0 \\ 4 \\ 16 \end{pmatrix}
		\]
		
		解线性方程组：
		\[
		\begin{pmatrix} 4 & 8 & 4 \\ 8 & 20 & 24 \\ 4 & 24 & 84 \end{pmatrix}
		\begin{pmatrix} x_1 \\ x_2 \\ x_3 \end{pmatrix}
		=
		\begin{pmatrix} 0 \\ 4 \\ 16 \end{pmatrix}
		\]
		
		\textbf{解得：}
		\[
		x = \begin{pmatrix} -2 \\ 1 \\ 0 \end{pmatrix}
		\]
		
		(\textit{注：代入验证，$4(-2)+8(1)=0$; $8(-2)+20(1)=4$; $4(-2)+24(1)=16$，成立。})
		
	\end{enumerate}
	
\end{document}